\addcontentsline{toc}{chapter}{Abstract}\vspace{-1cm}
%Border
\begin{tikzpicture}[remember picture, overlay]
  \draw[line width = 4pt] ($(current page.north west) + (0.75in,-0.75in)$) rectangle ($(current page.south east) + (-0.75in,0.75in)$);
\end{tikzpicture}

Modern warehouse operations face significant challenges in meeting the demands of rapidly expanding e-commerce and supply chain requirements, with traditional manual handling systems proving inadequate for achieving the speed, accuracy, and cost-effectiveness required in contemporary logistics environments. This research addresses these critical issues through the development of an intelligent autonomous delivery vehicle integrated with advanced computer vision capabilities, precision robotic manipulation systems, and cloud-based monitoring infrastructure using the NodeMCU ESP8266 platform.

The primary objective of this work is to develop a fully autonomous warehouse delivery system capable of intelligent object recognition, precise manipulation, and efficient navigation through complex warehouse environments. The system employs the LLaVA-4 model for real-time object classification integrated with servo-based robotic control systems for precise item handling operations.  \gls{ic}

The system development utilized Arduino IDE for embedded programming, OpenCV libraries for computer vision processing, and Firebase cloud platform for real-time data management and user interface development. The system's energy consumption analysis revealed 8.5 watts average power draw during active operation, representing 40\% improvement in energy efficiency compared to conventional automated warehouse systems. Navigation accuracy tests showed 98\% success rate with 45-second average completion time per task. \gls{ic}

Hardware implementation successfully validated simulation results through development of a functional prototype utilizing NodeMCU ESP8266 microcontroller, three servo motors for robotic arm articulation, L298N motor drivers for vehicle propulsion, and integrated camera systems for real-time object detection. Performance metrics from hardware testing closely matched simulation predictions with variations less than 5\% in execution times and accuracy measurements, demonstrating practical viability for real-world warehouse deployment.

\pagebreak